\chapter{Related Work and the Novelty Gap}
\label{ch:related-work-novelty-gap}

\section{Conformal prediction and adaptive calibration}
Conformal prediction and related distribution-free uncertainty methods provide one of the
few disciplined ways to wrap learned models with coverage-aware uncertainty without
assuming a perfectly specified error law
\cite{vovk2005algorithmic,lei2018distribution,romano2019conformalized,gibbs2021adaptive,angelopoulos2023gentle,barber2023beyond}.
Conformalized quantile regression and adaptive conformal variants extend that logic to
heteroskedastic and nonstationary settings, which makes them especially relevant once a
physical system moves through regime changes or degraded telemetry states
\cite{romano2019conformalized,gibbs2021adaptive,zaffran2022adaptive,stankeviciute2021conformal,xu2021dynamic,sesia2021comparison,bates2021distribution}.

\paragraph{Why this is insufficient for ORIUS.}
This family solves uncertainty quantification; it does not solve the release-boundary
question that follows from that uncertainty. Coverage guarantees do not by themselves decide
how the admissible action set should change, which fallback is allowed, or what evidence
must accompany release once the observation surface is unreliable. ORIUS therefore uses
calibration as an input to runtime safety, not as the full safety object.

\section{Runtime assurance and Simplex-style supervision}
Runtime assurance architectures, supervisory controllers, shielding, and runtime
verification supply the language of intervention, veto, fallback, and monitored execution
\cite{sha2001using,sha1998simplex,desai2019soter,bloem2015shield,bartocci2018introduction,havelund2004runtime,leucker2009brief,schierman2015run,bak2011runtime}.
That line of work is indispensable because it treats safe release as a separate problem from
nominal policy design.

\paragraph{Why this is insufficient for ORIUS.}
Most supervisory schemes enter after the nominal controller has already proposed an action
under a state estimate that is treated as sufficiently meaningful. ORIUS addresses an
earlier fault line: once observation quality degrades, the assurance layer must reason
jointly about reliability, widened state uncertainty, repaired action, fallback, and
certificate semantics. A supervisor without that observation-side logic is still missing the
central OASG object.

\section{Safety filters, shields, and robust safe control}
Control barrier functions, predictive safety filters, robust MPC, and tube-based control
translate model uncertainty into admissible control sets and corrective intervention rules
\cite{ames2019control,ames2014cbf,wabersich2021predictive,nguyen2016exponential,wang2017safety,fisac2019general,mayne2005robust,langson2004tubes,rawlings2017mpc,ben2009robust,camacho2013model}.
These methods are the closest classical relatives of ORIUS because they formalize the
difference between nominal control intent and safe control execution.

\paragraph{Why this is insufficient for ORIUS.}
This family usually assumes that the uncertainty set and safety object have already been
defined in a domain-specific way. ORIUS does not replace that work with a universal
controller. Instead, it provides a typed runtime contract that can host many repair
mechanisms while keeping reliability estimation, fallback semantics, and certificate
discipline explicit. The novelty is not the existence of safe control alone; it is the
runtime composition that makes safe control portable across domains.

\section{Anomaly, drift, and degraded-observation detection}
Fault diagnosis, change detection, cyber-physical security, and anomaly monitoring provide
the vocabulary for stale, delayed, reordered, dropped, or corrupted observations
\cite{rajkumar2010cps,derler2012modeling,baheti2011cyberphysical,pasqualetti2013attack,teixeira2015secure,basseville1993detection,page1954continuous,page1955test,hinkley1971inference,liu2008isolation,breunig2000lof}.
This literature is essential because degraded observation must first be recognized before it
can influence action release.

\paragraph{Why this is insufficient for ORIUS.}
Detection alone does not specify how action semantics should change after trust degrades. A
drift alarm can tell the system that something is wrong without saying how uncertainty
inflates, how the safe set tightens, what repair is permitted, or what certificate makes
the next step auditable. ORIUS uses detection as the opening move in a runtime safety
sequence, not as the full response.

\section{Domain-specific cyber-physical safety systems}
Energy systems, autonomous mobility stacks, industrial control, healthcare monitoring, and
other cyber-physical systems already contain rich safety practice in their own local forms
\cite{rajkumar2010cps,baheti2011cyberphysical,pasqualetti2013attack,teixeira2015secure,sha2001using,ames2019control}.
This practice matters because any claim of universality that ignores domain structure is
empty from the start.

\paragraph{Why this is insufficient for ORIUS.}
Domain-specific systems typically optimize within one application vocabulary and one local
evidence regime. They rarely expose a common contract that allows degraded-observation
replay, repaired action, fallback semantics, certificate validity, and latency to be
compared across domains. ORIUS addresses that missing cross-domain runtime object rather
than trying to erase domain structure.

\section{The novelty gap ORIUS addresses}
The literature becomes most powerful when these families are read together, but it also
reveals a persistent separation of levels. Uncertainty methods stop at calibrated sets.
Assurance methods stop at supervisory intervention. Safe-control methods stop at
plant-specific repair geometry. Detection methods stop at identifying degraded telemetry.
Domain systems stop at application-local safety practice. What is missing between them is a
single typed runtime composition that starts with observation reliability and ends with
repaired action, fallback, certificate semantics, and cross-domain replay under one common
contract.

That missing composition is ORIUS's novelty claim. The dissertation does not argue that any
one ORIUS component is unprecedented in isolation. It argues that the combination of
observation reliability, action repair, certificate semantics, and universal replay has not
previously been stabilized as one reusable runtime safety layer for physical AI. ORIUS is
therefore not ``better forecasting,'' ``better perception,'' or ``one more controller.'' It
is a release-boundary discipline for degraded-observation safety.

The closest plausible rival to ORIUS is not any single paper in these families. It is an
imagined composition of them: adaptive conformal intervals for uncertainty, a safety filter
or Simplex supervisor for intervention, anomaly detection for degraded telemetry, and
runtime verification for logging and policy enforcement. Even that composite stack still
stops short of ORIUS's central object. It lacks a typed release contract that binds
observed-state and true-state semantics, repaired action, fallback, certificate lifecycle,
and cross-domain replay under one explicit interface. That final integration step is the
novelty claim the dissertation defends.
